\chapter{Series}

A series adds up the terms of a sequence. Since infinitely many
additions cannot be carried out one after another, the value of a series is
\emph{defined} as the limit of its partial sums: series theory is sequence
theory applied to one particular sequence. This chapter covers the
convergence basics: the two sum formulas everyone must know (arithmetic and
geometric), the necessary condition for convergence with the harmonic series
as its famous counterexample, telescoping sums, and two light convergence
tests. The star of the chapter is the geometric series: the entire financial
mathematics of the next chapter (compound interest, present value,
annuities) is the geometric series in a suit.

\section{Core Rules}

\begin{rulebox}{Partial sums and convergence}
Given a sequence $(a_k)_{k\ge 1}$, its $n$-th \emph{partial sum} is
\[
  s_n \;=\; \sum_{k=1}^{n} a_k \;=\; a_1+a_2+\dots+a_n .
\]
The series $\sum_{k=1}^{\infty} a_k$ \emph{converges} if the sequence
$(s_n)$ of partial sums converges, and then
\[
  \sum_{k=1}^{\infty} a_k \;=\; \lim_{n\to\infty} s_n .
\]
If $(s_n)$ has no limit, the series \emph{diverges}. A series is thus
nothing but a sequence in disguise: every statement about series is a
statement about its sequence of partial sums.
\end{rulebox}

\begin{rulebox}{Arithmetic series (finite sum)}
For an arithmetic sequence $a_k=a+(k-1)d$ the sum of the first $n$ terms is
\[
  \sum_{k=1}^{n}\bigl(a+(k-1)d\bigr)
  \;=\; \frac{n}{2}\bigl(2a+(n-1)d\bigr)
  \;=\; n\cdot\frac{a_1+a_n}{2}
\],
$n$ times the average of the first and the last term. The special case
$a=d=1$ is Gauss' formula
\[
  \sum_{k=1}^{n} k \;=\; \frac{n(n+1)}{2} .
\]
An \emph{infinite} arithmetic series never
converges unless all terms are $0$: its terms do not tend to $0$.
\end{rulebox}

\begin{rulebox}{Geometric series, finite sum}
For a geometric sequence $a q^{k}$ with ratio $q\neq 1$,
\[
  \sum_{k=0}^{n-1} a q^{k}
  \;=\; a\,\frac{1-q^{n}}{1-q} \qquad (n \text{ terms}).
\]
For $q=1$ the sum is simply $na$.
\end{rulebox}

\begin{rulebox}{Geometric series, infinite sum}
If $|q|<1$, then $q^{n}\to 0$, so the partial sums converge and
\[
  \sum_{k=0}^{\infty} a q^{k} \;=\; \frac{a}{1-q},
  \qquad\text{and starting at } k=1:\quad
  \sum_{k=1}^{\infty} a q^{k} \;=\; \frac{aq}{1-q} .
\]
For $|q|\ge 1$ the series diverges. This is \textbf{the} key formula of the
whole book: in the next chapter every present value, annuity and perpetuity
formula of financial mathematics is obtained by choosing $q$ as a discount
factor in the geometric series.
\end{rulebox}

\begin{rulebox}{Necessary condition: the terms must vanish}
If $\sum_{k=1}^{\infty} a_k$ converges, then $a_k\to 0$. In contrapositive
form this is the everyday \emph{divergence test}:
\[
  a_k \not\to 0 \quad\Longrightarrow\quad \sum_{k=1}^{\infty} a_k
  \text{ diverges.}
\]
The converse is \textbf{false}: $a_k\to 0$ does \emph{not} guarantee
convergence. The standard counterexample is the harmonic series
\[
  \sum_{k=1}^{\infty}\frac{1}{k} \;=\; 1+\frac12+\frac13+\frac14+\dots,
\]
whose terms tend to $0$ while the partial sums grow beyond every bound.
\end{rulebox}

\begin{rulebox}{Telescoping sums}
If the terms can be written as consecutive differences, $a_k=b_k-b_{k+1}$,
then everything in the middle cancels:
\[
  \sum_{k=1}^{n}(b_k-b_{k+1}) \;=\; b_1-b_{n+1} .
\]
The series converges precisely when $(b_n)$ converges, with sum
$b_1-\lim_{n\to\infty} b_n$. A telescoping sum is
usually uncovered by partial fractions, e.g.
$\dfrac{1}{k(k+1)}=\dfrac{1}{k}-\dfrac{1}{k+1}$.
\end{rulebox}

\begin{rulebox}{Comparison test (simple form)}
Suppose $0\le a_k\le b_k$ for all $k$ (or from some index on).
\begin{itemize}[leftmargin=1.6em,itemsep=0.3ex]
  \item \textbf{Majorant}: if
        $\sum b_k$ converges, so does $\sum a_k$.
  \item \textbf{Minorant}: if
        $\sum a_k$ diverges, so does $\sum b_k$.
\end{itemize}
The usual benchmarks are the geometric series (convergent for $|q|<1$) and
the harmonic series (divergent). Example:
$\dfrac{1}{2^{k}+1}\le\Bigl(\dfrac12\Bigr)^{k}$, so
$\sum\dfrac{1}{2^{k}+1}$ converges.
\end{rulebox}

\begin{rulebox}{Ratio test (light form)}
For $a_k>0$ suppose the ratio of consecutive terms has a limit,
$\dfrac{a_{k+1}}{a_k}\to L$. Then:
\begin{itemize}[leftmargin=1.6em,itemsep=0.3ex]
  \item $L<1$: the series converges (it behaves like a geometric series
        with ratio $L$);
  \item $L>1$: the series diverges (the terms eventually grow);
  \item $L=1$: no conclusion: both $\sum 1/k$ (divergent) and
        $\sum 1/k^{2}$ (convergent) give $L=1$.
\end{itemize}
Example: for $a_k=k/2^{k}$ the ratio is
$\dfrac{k+1}{2k}\to\dfrac12<1$, so $\sum k/2^{k}$ converges. The test shines when powers or factorials appear.
\end{rulebox}

\begin{rulebox}{Linearity}
If $\sum a_k=A$ and $\sum b_k=B$ both converge and $\lambda,\mu\in\R$, then
\[
  \sum_{k=1}^{\infty}\bigl(\lambda a_k+\mu b_k\bigr)
  \;=\; \lambda A+\mu B .
\]
Splitting a series into pieces is only legitimate when \emph{each} piece
converges on its own.
\end{rulebox}

\section{Practice Problems}

Evaluate the sums in Problems 1--14. In Problems 15--19 decide convergence
or divergence with a justified test. Problem 20 previews the next chapter.

\begin{problems}
  \item $\displaystyle\sum_{k=1}^{100} k$
  \item $7+10+13+\dots+118$
  \item $\displaystyle\sum_{k=1}^{50}(2k-1)$ \quad (the first $50$ odd
        numbers)
  \item $\displaystyle\sum_{k=0}^{9} 2^{k}$
  \item $\displaystyle\sum_{k=1}^{6} 5\cdot 3^{k-1}$
  \item $\displaystyle\sum_{k=0}^{5}\Bigl(\frac12\Bigr)^{k}$
  \item $\displaystyle\sum_{k=0}^{\infty}\Bigl(\frac13\Bigr)^{k}$
  \item $\displaystyle\sum_{k=1}^{\infty}\Bigl(\frac25\Bigr)^{k}$
  \item $\displaystyle\sum_{k=0}^{\infty}\Bigl(-\frac12\Bigr)^{k}$
  \item Show that $0.999\ldots=0.\overline{9}$ equals exactly $1$.
  \item Write $2.353535\ldots=2.\overline{35}$ as a fraction.
  \item $\displaystyle\sum_{k=1}^{n}\frac{1}{k(k+1)}$ \quad (find a closed
        formula, then the value of the infinite series)
  \item $\displaystyle\sum_{k=1}^{99}\bigl(\sqrt{k+1}-\sqrt{k}\bigr)$
  \item $\displaystyle\sum_{k=1}^{\infty}\frac{1}{(2k-1)(2k+1)}$
  \item Does $\displaystyle\sum_{k=1}^{\infty}\frac{k}{k+1}$ converge?
  \item The terms of $\displaystyle\sum_{k=1}^{\infty}\frac{1}{k}$ satisfy
        $\frac{1}{k}\to 0$. Does the series converge? Justify carefully.
  \item Does $\displaystyle\sum_{k=1}^{\infty}\frac{1}{2^{k}+1}$ converge?
  \item Does $\displaystyle\sum_{k=1}^{\infty}\frac{1}{\sqrt{k}}$ converge?
  \item Does $\displaystyle\sum_{k=1}^{\infty}\frac{k}{2^{k}}$ converge?
  \item \emph{Perpetuity teaser.} A contract pays $100$ euros at the end of
        every year, forever. At an interest rate of $5\%$ per year, a
        payment $k$ years from now is worth $100\cdot(1/1.05)^{k}$ euros
        today. What is the total value of the contract today?
\end{problems}

\section{Solutions}

\begin{solutions}
  \item Arithmetic with $a_1=1$, $a_{100}=100$, $n=100$:
        \[
          \sum_{k=1}^{100} k=\frac{100\cdot 101}{2}=5050 .
        \]

  \item Arithmetic with $a=7$, $d=3$. The last term fixes $n$:
        $118=7+3(n-1)$ gives $3(n-1)=111$, so $n=38$. Hence
        \[
          7+10+\dots+118=\frac{38}{2}\,(7+118)=19\cdot 125=2375 .
        \]

  \item Arithmetic with $a_1=1$, $a_{50}=99$, $n=50$:
        \[
          \sum_{k=1}^{50}(2k-1)=50\cdot\frac{1+99}{2}=50\cdot 50=2500 .
        \]
        In general the first $n$ odd numbers sum to $n^{2}$.

  \item Geometric, $a=1$, $q=2$, $10$ terms:
        \[
          \sum_{k=0}^{9}2^{k}=\frac{1-2^{10}}{1-2}=\frac{1-1024}{-1}=1023 .
        \]

  \item Pull out the factor $5$; the remaining sum is geometric with $q=3$
        and $6$ terms:
        \[
          \sum_{k=1}^{6}5\cdot 3^{k-1}
          =5\sum_{j=0}^{5}3^{j}
          =5\cdot\frac{1-3^{6}}{1-3}
          =5\cdot\frac{-728}{-2}=5\cdot 364=1820 .
        \]

  \item Geometric, $q=\tfrac12$, $6$ terms:
        \[
          \sum_{k=0}^{5}\Bigl(\frac12\Bigr)^{k}
          =\frac{1-(1/2)^{6}}{1-1/2}
          =2\Bigl(1-\frac{1}{64}\Bigr)=\frac{63}{32} .
        \]

  \item $|q|=\tfrac13<1$, so the infinite sum formula applies:
        \[
          \sum_{k=0}^{\infty}\Bigl(\frac13\Bigr)^{k}
          =\frac{1}{1-\tfrac13}=\frac{3}{2} .
        \]

  \item The sum starts at $k=1$, so the first term is $\tfrac25$ and
        $q=\tfrac25$:
        \[
          \sum_{k=1}^{\infty}\Bigl(\frac25\Bigr)^{k}
          =\frac{2/5}{1-2/5}=\frac{2/5}{3/5}=\frac{2}{3} .
        \]

  \item The formula also covers negative ratios; here $q=-\tfrac12$ with
        $|q|<1$:
        \[
          \sum_{k=0}^{\infty}\Bigl(-\frac12\Bigr)^{k}
          =\frac{1}{1-(-\tfrac12)}=\frac{1}{3/2}=\frac{2}{3} .
        \]

  \item The decimal expansion is a geometric series with $q=\tfrac{1}{10}$:
        \[
          0.\overline{9}
          =\sum_{k=1}^{\infty}\frac{9}{10^{k}}
          =\frac{9\cdot\tfrac{1}{10}}{1-\tfrac{1}{10}}
          =\frac{9/10}{9/10}=1 .
        \]
        The equality is exact: $0.\overline{9}$ is not ``a little less
        than $1$'', it \emph{is} $1$.

  \item Split off the integer part; the repeating block $35$ has period
        $2$, so $q=\tfrac{1}{100}$:
        \[
          2.\overline{35}
          =2+\sum_{k=1}^{\infty}\frac{35}{100^{k}}
          =2+\frac{35\cdot\tfrac{1}{100}}{1-\tfrac{1}{100}}
          =2+\frac{35}{99}=\frac{233}{99} .
        \]

  \item Partial fractions: $\dfrac{1}{k(k+1)}=\dfrac{1}{k}-\dfrac{1}{k+1}$,
        so the sum telescopes:
        \[
          s_n=\sum_{k=1}^{n}\Bigl(\frac{1}{k}-\frac{1}{k+1}\Bigr)
          =1-\frac{1}{n+1} .
        \]
        Since $\dfrac{1}{n+1}\to 0$,
        \[
          \sum_{k=1}^{\infty}\frac{1}{k(k+1)}=\lim_{n\to\infty} s_n=1 .
        \]

  \item Telescoping with $b_k=-\sqrt{k}$ (or simply: all middle terms
        cancel):
        \[
          \sum_{k=1}^{99}\bigl(\sqrt{k+1}-\sqrt{k}\bigr)
          =\sqrt{100}-\sqrt{1}=10-1=9 .
        \]

  \item Partial fractions:
        \[
          \frac{1}{(2k-1)(2k+1)}
          =\frac12\Bigl(\frac{1}{2k-1}-\frac{1}{2k+1}\Bigr).
        \]
        The partial sums telescope:
        \[
          s_n=\frac12\Bigl(1-\frac{1}{2n+1}\Bigr)
          \;\longrightarrow\;\frac12 ,
        \]
        so the series converges with sum $\dfrac12$.

  \item No. The terms satisfy
        \[
          a_k=\frac{k}{k+1}\;\longrightarrow\;1\neq 0,
        \]
        so the necessary condition fails and the series diverges. (Each
        late term adds almost $1$, so the partial sums grow roughly like
        $n$.)

  \item No: this is the harmonic series, the standard warning that
        $a_k\to 0$ is necessary but \emph{not sufficient}. Group the terms
        in blocks of doubling length:
        \[
          1+\frac12
          +\underbrace{\Bigl(\frac13+\frac14\Bigr)}_{\ge\,2\cdot\frac14=\frac12}
          +\underbrace{\Bigl(\frac15+\frac16+\frac17+\frac18\Bigr)}_{\ge\,4\cdot\frac18=\frac12}
          +\dots
        \]
        Each block contributes at least $\tfrac12$, so
        $s_{2^{m}}\ge 1+\tfrac{m}{2}$. The partial sums exceed every bound,
        hence the series diverges.

  \item Yes, by comparison with a geometric majorant: for all $k\ge 1$,
        \[
          0<\frac{1}{2^{k}+1}<\Bigl(\frac12\Bigr)^{k},
          \qquad
          \sum_{k=1}^{\infty}\Bigl(\frac12\Bigr)^{k}
          =\frac{1/2}{1-1/2}=1<\infty ,
        \]
        so $\displaystyle\sum_{k=1}^{\infty}\frac{1}{2^{k}+1}$ converges.

  \item No. For $k\ge 1$ we have $\sqrt{k}\le k$, hence
        \[
          \frac{1}{\sqrt{k}}\ \ge\ \frac{1}{k} .
        \]
        The harmonic series diverges, so by the minorant criterion
        $\displaystyle\sum_{k=1}^{\infty}\frac{1}{\sqrt{k}}$ diverges as
        well.

  \item Yes, by the ratio test with $a_k=k/2^{k}$:
        \[
          \frac{a_{k+1}}{a_k}
          =\frac{k+1}{2^{k+1}}\cdot\frac{2^{k}}{k}
          =\frac{k+1}{2k}
          \;\longrightarrow\;\frac12<1 ,
        \]
        so the series converges. (Its value happens to be $2$, but the
        test alone does not reveal that.)

  \item The total present value is a geometric series with
        $q=\dfrac{1}{1.05}$, starting at $k=1$:
        \[
          \mathrm{PV}
          =\sum_{k=1}^{\infty}100\,\Bigl(\frac{1}{1.05}\Bigr)^{k}
          =100\cdot\frac{\tfrac{1}{1.05}}{1-\tfrac{1}{1.05}}
          =100\cdot\frac{\tfrac{1}{1.05}}{\tfrac{0.05}{1.05}}
          =\frac{100}{0.05}=2000 .
        \]
        An infinite stream of payments is worth a finite amount:
        $2000$ euros. The rule of thumb
        $\mathrm{PV}=\text{payment}/\text{interest rate}$ is exactly the
        infinite geometric sum, the next chapter builds all of financial
        mathematics on this observation.
\end{solutions}

\section{Additional Problems}

\begin{problems}
  \item $\displaystyle\sum_{k=1}^{200}(3k+2)$
  \item $100+96+92+\dots+4$
  \item $\displaystyle\sum_{k=0}^{7}3\cdot(-2)^{k}$
  \item $\displaystyle\sum_{k=2}^{\infty}\Bigl(\frac14\Bigr)^{k}$
  \item Write $0.727272\ldots=0.\overline{72}$ as a fraction.
  \item $\displaystyle\sum_{k=1}^{\infty}\frac{1}{(k+1)(k+2)}$
  \item Does $\displaystyle\sum_{k=1}^{\infty}\frac{2k+1}{3k+5}$ converge?
  \item Show that $\displaystyle\sum_{k=1}^{\infty}\frac{1}{k^{2}}$
        converges, by comparing with the telescoping series of
        Problem~12 of the practice set
        \emph{(hint: $\tfrac{1}{k^{2}}\le\tfrac{2}{k(k+1)}$)}.
  \item Does $\displaystyle\sum_{k=1}^{\infty}\frac{2^{k}}{k!}$ converge?
  \item $\displaystyle\sum_{k=0}^{\infty}
        \Bigl(\Bigl(\frac12\Bigr)^{k}+2\cdot\Bigl(\frac13\Bigr)^{k}\Bigr)$
\end{problems}

\section{Additional Solutions}

\begin{solutions}
  \item Split by linearity of finite sums and use Gauss' formula:
        \[
          \sum_{k=1}^{200}(3k+2)
          =3\sum_{k=1}^{200}k+\sum_{k=1}^{200}2
          =3\cdot\frac{200\cdot 201}{2}+400
          =3\cdot 20100+400=60700 .
        \]

  \item Arithmetic with $a=100$, $d=-4$. The last term fixes $n$:
        $4=100-4(n-1)$ gives $4(n-1)=96$, so $n=25$. Hence
        \[
          100+96+\dots+4=\frac{25}{2}\,(100+4)=25\cdot 52=1300 .
        \]

  \item Geometric with $q=-2$, $8$ terms:
        \[
          \sum_{k=0}^{7}3\cdot(-2)^{k}
          =3\cdot\frac{1-(-2)^{8}}{1-(-2)}
          =3\cdot\frac{1-256}{3}=-255 .
        \]

  \item Geometric with $|q|=\tfrac14<1$; the first term is $(1/4)^{2}$:
        \[
          \sum_{k=2}^{\infty}\Bigl(\frac14\Bigr)^{k}
          =\frac{(1/4)^{2}}{1-1/4}
          =\frac{1/16}{3/4}=\frac{1}{12} .
        \]
        (Check: full sum $\tfrac{4}{3}$ minus the terms $1+\tfrac14$
        indeed leaves $\tfrac{4}{3}-\tfrac{5}{4}=\tfrac{1}{12}$.)

  \item Repeating block $72$ of period $2$, so $q=\tfrac{1}{100}$:
        \[
          0.\overline{72}
          =\sum_{k=1}^{\infty}\frac{72}{100^{k}}
          =\frac{72\cdot\tfrac{1}{100}}{1-\tfrac{1}{100}}
          =\frac{72}{99}=\frac{8}{11} .
        \]

  \item Partial fractions:
        $\dfrac{1}{(k+1)(k+2)}=\dfrac{1}{k+1}-\dfrac{1}{k+2}$, so
        \[
          s_n=\frac12-\frac{1}{n+2}\;\longrightarrow\;\frac12 ,
        \]
        and the series converges with sum $\dfrac12$.

  \item No. The terms satisfy
        \[
          \frac{2k+1}{3k+5}
          =\frac{2+\tfrac{1}{k}}{3+\tfrac{5}{k}}
          \;\longrightarrow\;\frac{2}{3}\neq 0 ,
        \]
        so the necessary condition fails and the series diverges.

  \item For $k\ge 1$ we have $k+1\le 2k$, hence $k(k+1)\le 2k^{2}$ and
        therefore
        \[
          \frac{1}{k^{2}}\ \le\ \frac{2}{k(k+1)} .
        \]
        The majorant is twice the telescoping series with sum $1$, so
        $\sum_{k=1}^{\infty}\tfrac{2}{k(k+1)}=2$ converges. By the
        comparison test, $\displaystyle\sum_{k=1}^{\infty}\frac{1}{k^{2}}$
        converges (its partial sums are increasing and bounded by $2$).
        Its exact value, $\pi^{2}/6$, is a famous result of Euler: far
        beyond what the test tells us.

  \item Yes, by the ratio test with $a_k=2^{k}/k!$:
        \[
          \frac{a_{k+1}}{a_k}
          =\frac{2^{k+1}}{(k+1)!}\cdot\frac{k!}{2^{k}}
          =\frac{2}{k+1}
          \;\longrightarrow\;0<1 ,
        \]
        so the series converges. (Its value is $\mathrm{e}^{2}-1$, as the
        exponential series will show later.)

  \item Both pieces are convergent geometric series, so linearity applies:
        \[
          \sum_{k=0}^{\infty}\Bigl(\frac12\Bigr)^{k}
          =\frac{1}{1-\tfrac12}=2 ,
          \qquad
          2\sum_{k=0}^{\infty}\Bigl(\frac13\Bigr)^{k}
          =2\cdot\frac{1}{1-\tfrac13}=2\cdot\frac32=3 ,
        \]
        hence the total is $2+3=5$.
\end{solutions}
